\section{Local passage tools at a simple contact}
\label{sec:simple-contact-tools}

We treat a contact whose linear coefficient is negative, with estimates
uniform as that coefficient tends to zero. This uniformity allows the
same tools to be used at the quadratic contact in
Section~\ref{sec:positive-constant}.
Throughout this section the field is
\begin{equation}\label{eq:simple-family}
 F=y(1+bz-y)+\varepsilon^2(Cz^2-Az+A^2c),\qquad
 G=zy+\varepsilon^3A^3d.
\end{equation}
The constants in estimates are independent of the positive small
parameters.  The letter $C$ in \eqref{eq:simple-family} always denotes
the quadratic coefficient, not a generic constant.  All derivatives
in the estimates below are taken in phase or section variables, with
the field parameters fixed.

To bound the zeros of the complete displacement, we compare a local
contact passage with its complete complementary transition. We first
estimate the negative and positive transitions in regular section
coordinates, then construct the local saddle passages. These common
estimates will be used in both contact-scale regimes below.

The return maps below are defined using the sections of
Lemma~\ref{lem:sections}.  In this section $u$ denotes the physical
ordinate on its outgoing section $z=z_c>0$, where $F>0$; it is not the
coordinate on the upper section.  We put the negative section at
$z=-q_g$, $q_g>0$.
Every cycle being counted meets the positive section exactly once and
has, in the contact neighborhood, $y>0$, $F>0$ and
$y(-q_g)\le K_0\varepsilon^3$.  We decrease that neighborhood so that
$1/2<1+bz-y<2$ and $F_y>1/2$ there.

\subsection{The negative transition}

The negative transition contracts strongly enough to absorb the shrinking
transverse scale in both regimes. Write
\begin{equation}\label{eq:simple-tail-parameters}
 c=t^2\kappa,\quad d=t^3\eta,\quad \varepsilon=te,
 \qquad 0<e\le1,\quad |\kappa|,|\eta|\le K,
 \quad n=At^2.
\end{equation}
We shall take $t$ small; $e$ need not be bounded away from zero.
Set $p=\varepsilon^2P$ and $R=G/F$.

The central contact passage uses the finite-order saddle-node product of
\cite[Theorem~3 and Section~2.8]{DIR2002}, at pp.~6 and 11 of
the authors' public manuscript. We apply it in the central contact chart
of Appendix~\ref{sec:simple-contact}.
Its planar critical germ has one zero and one negative eigenvalue;
the center reduction has a nonzero quadratic term. These data match
the multiplicity-two translated critical field there; parameters
form a finite-dimensional smooth unfolding. The differentiability
order is prescribed before choosing its smaller parameter window.
It supplies a finite-smooth orbital product. The fixed-gap product used
away from the central saddle-node is proved in
Section~\ref{sec:normal-product}. Neither construction uses analytic
linearization across degeneration or a positive-endpoint cyclicity bound.

\begin{lemma}\label{lem:negative-transition}
Assume $0<A\le A_+<\infty$ and either
\begin{enumerate}
\item $C$ belongs to a fixed compact subinterval of $(0,\infty)$; or
\item $A\ge A_0>0$, $C$ is bounded, and $P(z)\ge p_0>0$
on the fixed negative interval preceding the contact neighborhood.
\end{enumerate}
In the second case choose the fixed negative cut $q_g$ sufficiently
small in terms of these bounds. With the scales in \eqref{eq:simple-tail-parameters}, fix $L$
sufficiently large and then take $t$ sufficiently small. Every positive,
right-going transition from the prescribed physical source, with
$y(-q_g)\le K_0\varepsilon^3$, satisfies on its passage from
$z=-q_g$ to $z=-Ln$
\begin{equation}\label{eq:simple-negative-barrier}
 0<y<p,\qquad p\le F\le3p,\qquad
 P(-x)\asymp x(x+A)\quad(Ln\le x\le q_g).
\end{equation}
If $J_-(u)=y(-Ln;u)$ includes the preceding fast part of the orbit, then
\begin{equation}\label{eq:simple-negative-jets}
 0<(J_-)_u\le K_1\varepsilon e^{-\gamma T_A/\varepsilon^2},
 \qquad 0\le \frac{(J_-)_{uu}}{(J_-)_u}\le K_1\varepsilon^{-2},
 \qquad T_A=\log\frac{q_g+A}{A(1+Lt^2)}.
\end{equation}
Here $T_A\ge t_0+t_1\log(A_+/A)$ for some $t_0,t_1>0$.
The transverse ordinate scale is $n^2e^2=A^2\varepsilon^2t^2$.
In particular,
\begin{equation}\label{eq:simple-scale-cancellation}
 \frac{(J_-)_u}{A^2\varepsilon^2t^2}\longrightarrow0
\end{equation}
uniformly in $A,t,e$.
\end{lemma}

\begin{proof}
If $C$ belongs to a fixed compact subinterval of $(0,\infty)$,
$P(-x)\ge Cx^2+Ax/2$ for $x\ge LAt^2$ after taking $L\ge2K$.
If instead $A\ge A_0>0$, $C$ is bounded, and $P$ has a fixed positive
lower bound on the preceding negative interval, choose $q_g$ with
$|C|q_g\le A_0/8$ and
$K/L\le1/8$.  Then
$3Ax/4\le P(-x)\le5Ax/4$, which gives the same comparability in
\eqref{eq:simple-negative-barrier}.  At the graph $y=p(z)$,
\[
 G-p_zF\le\varepsilon^2xP(-x)
 \left(-1+\frac{3e^2}{L}+6|C|\varepsilon^2
                  +\frac{2eK}{L^2}\right)<0.
\]
The incoming height is below $p(-q_g)$ for small $\varepsilon$.
This proves the barrier by first contact.

For later use we record the exact derivatives.  With $Q=1+bz$ and
$\Delta=\varepsilon^3A^3t^3\eta$,
\begin{align}
 R_y&=\frac{z(y^2+p)-\Delta(Q-2y)}{F^2},\label{eq:simple-Ry}\\
 R_{yy}&=\frac{2\{z[y^3+(3y-Q)p]
              +\Delta[p+Q^2-3Qy+3y^2]\}}{F^3}.
              \label{eq:simple-Ryy}
\end{align}
On the small negative tail $|\Delta|/(xp)\le2eK/L^2$.
The numerator $z[y^3+(3y-Q)p]$ is positive and comparable to $xp$.
Increasing $L$, and decreasing the contact neighborhood, therefore gives
\begin{equation}\label{eq:simple-tail-density}
 R_y\le-\frac{\gamma}{\varepsilon^2(x+A)},\qquad
 0<R_{yy}\le\frac{K_2}{\varepsilon^4x(x+A)^2}.
\end{equation}

The complete negative transition also includes an entrance whose
abscissa moves with $u$.
Follow the descending part of the fast orbit to
$y_e=K_e\varepsilon^3$.  Its abscissa $\Lambda(u)$ lies in a fixed
negative interval, with $0<k_1\le\Lambda_u\le k_2$ and
$|\Lambda_{uu}|\le k_2$.  Indeed, on the descending part,
\[
 \frac{dz}{dy}=
 \frac{1+bz-y+\varepsilon^2P(z)/y}{z+\Delta/y},
 \qquad
 \int_{y_e}^{y_h}\left(1+\frac{\varepsilon^2}{y}\right)dy\le K_3;
\]
the denominator stays away from zero and the first two variational
equations give these bounds.  On the subsequent fixed negative interval,
$R=-\Theta(\varepsilon)$ at the entrance,
$R_y\le-k_3\varepsilon^{-2}$ and
$|R_{yy}|\le K_3\varepsilon^{-4}$.  The derivative and its logarithmic
derivative, at a fixed terminal abscissa, are exactly
\begin{align}
 y_u(z)&=-R(\Lambda,y_e)\Lambda_u
                   \exp\!\int_\Lambda^zR_y,\label{eq:simple-moving-first}\\
 \partial_u\log y_u(z)&=
 \frac{\Lambda_{uu}}{\Lambda_u}
 +\left(\frac{R_z}{R}-R_y\right)(\Lambda,y_e)\Lambda_u
 +\int_\Lambda^zR_{yy}y_u.\label{eq:simple-moving-second}
\end{align}
Here $R_z/R=O(1)$ at the entrance.  The last integral on the fixed
interval has absolute value at most
$K_4\varepsilon^{-4}\int_0^{O(1)}\varepsilon
e^{-k_3s/\varepsilon^2}ds=O(\varepsilon^{-1})$.
The term $-R_y\Lambda_u$ is at least $k_4\varepsilon^{-2}$.
Thus the map to $-q_g$ is increasing and convex, with logarithmic
derivative bounded by $K_4\varepsilon^{-2}$.  This argument does not
require $R_{yy}$ to be positive on the remote interval.

On the tail, \eqref{eq:simple-tail-density} preserves convexity and
bounds $(J_-)_u$. Its contribution to the second variation reduces,
with $N=\gamma/\varepsilon^2$, to
\begin{equation}\label{eq:simple-tail-integral}
 \frac{K_5\varepsilon^{-3}}{(q_g+A)^N}
       \int_{LAt^2}^{q_g}\frac{(x+A)^{N-2}}x\,dx.
\end{equation}
For $A\le q_g/4$, split at $x=A$.  The upper integral, using
$x^{-1}\le2/(x+A)$, is $O(\varepsilon^{-1})$; the lower one is bounded
by
\[
 K_5\varepsilon^{-3}|\log(Lt^2)|A^{-2}
       \left(\frac{2A}{q_g+A}\right)^N=o(1).
\]
For large $N$ the $A$-dependent factor is increasing up to $q_g/4$,
where its exponential ratio is $2/5$.
For $q_g/4\le A\le A_+$ split instead at $x=q_g/2$.
The upper integral is again $O(\varepsilon^{-1})$, and the lower one
is bounded by
$K_6\varepsilon^{-3}(1+|\log\varepsilon|)\theta^N=o(1)$,
where $\theta=(q_g/2+A_+)/(q_g+A_+)<1$.
We used $t\ge\varepsilon$.  In the signed-$C$ case one may equivalently
use $P(-x)\asymp Ax$, split at $q_g/2$, and obtain the same estimate.
This proves \eqref{eq:simple-negative-jets}.

Choose $t$ so that $T_A(A_+,t)>t_0>0$.  Differentiation with respect
to $\log A$ gives
$\partial_{\log A}T_A=-q_g/(q_g+A)$, proving its asserted lower
bound.  Finally the left side of \eqref{eq:simple-scale-cancellation}
is at most
\[
 K_7\varepsilon^{-3}A^{-2}
       \left(\frac{A(1+Lt^2)}{q_g+A}\right)^{\gamma/\varepsilon^2}.
\]
For $N>2+2A_+/q_g$, the function
$A^{-2}[A/(q_g+A)]^N$ is increasing.  Its maximum is at $A_+$,
where the displayed exponential ratio is uniformly below one.
The exponential therefore absorbs the shrinking ordinate scale, with
no restriction on the relative rates of $A,t,e$.
\end{proof}

\subsection{Positive transitions and local passage maps}

On the positive transition, a relative transverse coordinate avoids a
second-derivative loss over a long flight. The following variational
estimate controls that coordinate by its total change.

\begin{lemma}\label{lem:normal-distortion}
Suppose $v_x=\alpha(x,v)v$ on a forward interval
$x_0\le x\le x_1$, with $x_1>x_0$ and $v(x_0,a)=a>0$, and suppose the
selected solutions are monotone in $x$ and remain below $v_0$.  If
\[
 |v\alpha_v|+|v^2\alpha_{vv}|\le K v|\alpha|,
\]
then the endpoint map $\Phi$ satisfies
$a\Phi'/\Phi=1+O(v_0)$ and $a\Phi''/\Phi'=O(v_0)$.
If $\alpha\le0$, the stronger estimates
\begin{equation}\label{eq:simple-contracting-distortion}
 \Phi'>0,\qquad \Phi'\le K_1,\qquad
 |\Phi''/\Phi'|\le K_1
\end{equation}
hold, independently of its flight time or of $\Phi'(0)$.
\end{lemma}

\begin{proof}
Put $E=av_a/v$ and $B=a\partial_a\log E$.  Exact differentiation gives
\[
 (\log E)_x=v\alpha_v,\qquad
 B_x=E(v\alpha_v+v^2\alpha_{vv}),\qquad E(x_0)=1, B(x_0)=0.
\]
Since $\int v|\alpha|dx=|v(x_1)-a|\le v_0$,
$|\log E|\le Kv_0$ and $|B|\le Ke^{Kv_0}v_0$.
The identity $a\Phi''/\Phi'=E-1+B$ proves the first assertion.
If $\alpha\le0$, replace $v_0$ by $a$ and use
$\Phi'=E\Phi/a\le E$.
\end{proof}

The next lemma converts this relative distortion bound into physical
section derivatives. Its input is a reference outgoing orbit of the
same field, with quadratic height bounds and positive speed. We will
construct that orbit separately for the bounded and small saddle ratios.

\begin{lemma}\label{lem:positive-transition}
Put $z=nU$, $y=n^2V$, where $n>0$ and either
$n=A\varepsilon^2$ or $n=Ac$.
Consider the scaled equations
\begin{equation}\label{eq:simple-positive-scaled}
 U'=V(1+bnU-n^2V)+q(U),\qquad V'=UV+\ell_1,
\end{equation}
where $q(U)=k-U+C\varepsilon^2U^2$, with bounded $k$, or
$q(U)=e^2(1-U)+C\varepsilon^2U^2$.
Let $\mathcal F(U,V)$ denote the right-hand side of the first
equation. Suppose a reference outgoing orbit of the same system, in
particular with the same value of $\ell_1$, is written as $V=h(U)$ and
reaches a fixed $U_0>0$ with a uniform positive speed and has the bounds
\begin{equation}\label{eq:simple-positive-reference}
 k_0(1+U^2)\le h(U)\le K_0(1+U^2),\qquad
 \mathcal F(U,h)\ge k_0h,
\end{equation}
up to a fixed small physical section $z=z_c$.
For $|\ell_1|$ sufficiently small, a fixed small relative neighborhood
of this reference reaches the same section.  If its initial transverse
coordinate $a$ is oriented so that
$V/h-1=j_0a+O(a^2)$, with $j_0$ a positive unit and bounded first
three derivatives, its physical transition $u=C_+(a)$ satisfies
\begin{equation}\label{eq:simple-positive-jets}
 0<C_+'\le K_1 n^{3/2},\qquad
 \left|\frac{C_+''}{C_+'}\right|\le K_1.
\end{equation}
The conclusions hold for bounded signed $C$.  The entire relative
neighborhood lies in the fixed outgoing section near the limiting fast
orbit, after decreasing the parameters.
The reference graph and the constants in
\eqref{eq:simple-positive-reference} are required only on the compact
$U$--interval ending at the fixed section $z=z_c$.  In each application
below that graph, its two barriers, and the positive-speed collar are
constructed first on this interval; the parameter neighborhood is then
shrunk so that the same collar is valid for all parameters under
consideration.  No assertion is made for $U$ outside this interval.
\end{lemma}

\begin{proof}
Write $V=h(1+v)$, $F_0=\mathcal F(U,h)$,
$F_v=\mathcal F(U,h(1+v))$ and
$B_v=1+bnU-n^2h(2+v)$.  Subtracting the two orbit equations gives
\begin{equation}\label{eq:simple-relative-equation}
 v_U=\alpha(U,v)v,\qquad
 \alpha=-\frac{(1+v)UhB_v+(\ell_1/h)F_0
                      +\ell_1(1+v)B_v}{F_vF_0}.
\end{equation}
On the physical neighborhood, $B_v$ is a positive unit. By
\eqref{eq:simple-positive-reference}, $F_v/F_0=1+O(v)$ and
$h/F_0$ is bounded.  The two terms involving $\ell_1$ are uniformly
smaller than the first numerator term, since $U\ge U_0>0$ and
$h\ge k_0(1+U^2)$.  Thus $\alpha<0$ and direct differentiation gives
\[
 |v\alpha_v|+|v^2\alpha_{vv}|\le K_2|v||\alpha|.
\]
Apply Lemma~\ref{lem:normal-distortion} to $v$ or $-v$.
The relative deviation contracts, so it cannot leave the chosen
neighborhood; the second estimate in \eqref{eq:simple-positive-jets}
follows, including the initial smooth change of section.

To recover the size of $C_+'$, use \eqref{eq:simple-Ry} along the
reference neighborhood. The physical constant and forcing terms have
sizes $O(n^2)$ and $O(n^3)$, and the term proportional to $-z^2$ has
a favorable sign.  Since $F\ge k_1y$ and $y\ge k_1(n^2+z^2)$,
\[
 R_y\le K_3z+
 K_3\frac{\varepsilon^2z^3+n^2z+n^3}{(n^2+z^2)^2}.
\]
Its integral to $z_c$ is at most
$K_4+K_4\varepsilon^2\log(1/n)$.  The physical transverse derivative
at the initial scaled section is $O(n^2)$.  Hence
$C_+'\le K_5n^{2-K_4\varepsilon^2}\le K_5n^{3/2}$, where the last
inequality uses one common choice $\varepsilon_0$ with
$K_4\varepsilon_0^2\le1/2$ and then the common shrinking $n<1$.
Finally $k_2z^2\le y\le K_6(n^2+z^2)$ on fixed positive intervals.
In the joint limit $n,\varepsilon\to0$, any limiting graph therefore solves
$y_z=z/(1+bz-y)$ and has $y(0)=0$.
Uniqueness for this regular equation identifies it with the selected
fast orbit.  The bounds give compactness on every $[z_0,z_c]$ and then
convergence as $z_0\downarrow0$, uniformly over the relative
neighborhood.  This proves the last assertion.
\end{proof}

At the central saddle-node, the exact product separates the scalar
flight time from normal contraction. It also reduces the fixed-point
count to the sign of one derivative involving the complement.

\begin{lemma}\label{lem:saddle-node-passage}
Let $X_p$ be a $C^\infty$ local family, with $p$ in a compact
parameter germ, and suppose the phase germ of $X_{p_0}$ is real analytic
and has an ordinary multiplicity-two saddle-node with a negative nonzero
eigenvalue and a positive quadratic center coefficient.  For every
integer $k<\infty$, after decreasing the
phase--parameter neighborhood, there are parameter-dependent $C^k$
orbital coordinates
\[
 \dot s=W(s,p),\qquad \dot w=-w,\qquad
 W=u(s,p)((s-\sigma(p))^2+\mu(p)),\quad u>0.
\]
The coordinates, their inverses, and the regular changes to fixed
transverse sections have bounded derivatives through every order fixed
before that shrinking.  No $C^\infty$ or analytic equivalence on one
common phase--parameter neighborhood is asserted.
Let $D(s)=w_0\exp(-\int_s^{s_+}W^{-1})$ be a finite passage from
$w=w_0>0$ to $s=s_+$.  Then $WD_s=D$.  On a sufficiently small
incoming interval,
\begin{equation}\label{eq:simple-saddle-node-smallness}
 D(s)\le w_0\exp\left(-\frac{k_0}
                         {|s-\sigma|+\sqrt{|\mu|}}\right).
\end{equation}
If an increasing return transition $H$ is defined on an interval of
$a>0$, and on one subinterval containing every solution of $D(H(a))=a$
one has $W(H)>0$ and
\begin{equation}\label{eq:simple-pfaff-hypothesis}
 1-W_s(H(a))+aH''(a)/H'(a)>0,
\end{equation}
then the equation has total isolated zero multiplicity at most two
on that actual finite-passage subinterval.
\end{lemma}

\begin{proof}
The finite-order orbital product form is the multiplicity-two
specialization of \cite[Section~2.1, equation~(2.2), Theorem~3, and
Section~2.8]{DIR2002}.  More precisely, the source treats a
$C^\infty$ unfolding of a real-analytic critical saddle-node germ.
Its multiplicity is denoted by $\mu+1$, so the present case has
$\mu=1$ and, after a possible reversal of time, its normal form is
\[
 \dot z=b_0(p)+\frac{z^2}{1+\alpha(p)z},\qquad \dot w=-w,
\]
because its polynomial $P(p,z)$ has degree at most $\mu-1=0$ in the
multiplicity-two case.  This is an exact product, not a formal jet.
Moreover
\[
 b_0+\frac{z^2}{1+\alpha z}
 =\frac{(z-\sigma)^2+\widehat\mu}{1+\alpha z},\qquad
 \sigma=-\frac{\alpha b_0}{2},\qquad
 \widehat\mu=b_0-\frac{\alpha^2b_0^2}{4}.
\]
Shrinking the phase neighborhood makes $(1+\alpha z)^{-1}$ a positive
unit, which gives the asserted $W$ after renaming $z$ and
$\widehat\mu$.  The source states finite smoothness on a
phase--parameter neighborhood and that decreasing this neighborhood
permits a higher prescribed finite order; Section~2.8 treats a
multidimensional unfolding.  We therefore include all compact passive
parameters in $p$ before choosing $k$ and shrinking.  Only the resulting
$C^k$ coordinates and regular $C^k$ section changes are used here; the
source's analyticity assertion away from the critical stable manifold
is not used.

In these canonical coordinates the second equation integrates exactly.
Along a passage on which $W>0$,
\[
 \frac{dw}{ds}=-\frac{w}{W(s,p)},\qquad
 D(s)=w_0\exp\left(-\int_s^{s_+}\frac{dx}{W(x,p)}\right),
 \qquad WD_s=D.
\]
Write $s_+=\sigma+r$ and
$m=\max(|s-\sigma|,\sqrt{|\mu|})\le r/2$.
Every finite passage under consideration contains the interval
$[\sigma+m,\sigma+r]$ and has
\[
 \int_s^{s_+}\frac{dx}{W(x)}\ge
 \int_m^r\frac{dx}{u_1(x^2+|\mu|)}\ge\frac1{4u_1m}.
\]
If $\mu\le0$, its starting point is necessarily to the right of the
rightmost root; no root is crossed.  This proves
\eqref{eq:simple-saddle-node-smallness}, including the passage from
$s>\sigma$ at $\mu=0$.
For $L(a)=\log D(H(a))-\log a$ and
$Z(a)=aH'(a)-W(H(a))$, exact differentiation gives
\[
 L'=\frac{Z}{aW(H)},\qquad
 Z'=H'\left(1-W_s(H)+aH''/H'\right)>0.
\]
Thus $L'$ has at most one zero. A fixed point with $L'\ne0$ is
simple; if $L=L'=0$, then $L''=Z'/[aW(H)]>0$, so it is exactly
quadratic and the only tangent zero. Since the critical point is a
strict minimum, the sum of isolated fixed-point multiplicities is at
most two on this actual finite-passage interval.
\end{proof}

When the saddle ratio tends to zero, we instead need product coordinates
whose phase derivatives remain bounded in that limit. The following
lemma constructs them and gives the corresponding weighted curvature
criterion for the complement.

\begin{lemma}\label{lem:small-ratio-saddle}
For the system
\begin{equation}\label{eq:simple-small-ratio-system}
 X'=e^2\{Y(1+bnX-n^2e^2Y)+1-X+CcX^2\},\qquad
 Y'=XY+e\eta,
\end{equation}
near $X=1,Y=0$, with $n=Ac$, $0<A\le A_+$,
$c,e,|\eta|$ small and bounded $C$, there are
uniform $C^j$ orbital coordinates, for any fixed $j$, in which
\begin{equation}\label{eq:simple-exact-power}
 \dot\xi=-r\xi,\quad \dot\omega=\omega,\qquad
 r\asymp e^2,\qquad D(s)=\xi_0(s/\omega_0)^r.
\end{equation}
Here $D$ is the exact passage between fixed positive sections
$\xi=\xi_0$ and $\omega=\omega_0$.
If an increasing complete complementary transition $H$ has
\begin{equation}\label{eq:simple-small-ratio-curvature}
 r|H''/H'|<\frac{1-r_0}{2\xi_0},\qquad 0<r\le r_0<1,
\end{equation}
then $D(H(a))=a$ has total isolated zero multiplicity at most two on
each actual finite-passage interval where these maps are defined.
\end{lemma}

\begin{proof}
It suffices to establish the assertion for $j\ge4$. Fix an integer
$K\ge2j+11$ before shrinking the parameters, and choose nested intervals
$I_0\Subset I_1\Subset I_2\Subset(0,\infty)$ about $X=1$.
The interval $I_0$ contains the saddle and the local sections to be used
below. Write
\[
 \mathcal F(X,Y)=Y(1+bnX-n^2e^2Y)+1-X+CcX^2,
 \qquad w=Y+\frac{e\eta}{X}.
\]
In the original time the translated equations are exactly
\[
 X'=e^2\mathcal F(X,w-e\eta/X),\qquad
 w'=Xw-\frac{e^3\eta}{X^2}\mathcal F(X,w-e\eta/X).
\]
Thus, after reversing time, they have the form
\begin{equation}\label{eq:small-ratio-translated}
 \dot X=e^2\mathcal A(X,w),\qquad
 \dot w=-Xw+q\mathcal R(X,w),\qquad q=e^3\eta,
\end{equation}
where $\mathcal A=-\mathcal F(X,w-e\eta/X)$ and
$\mathcal R=X^{-2}\mathcal F(X,w-e\eta/X)$ have common finite-order
phase bounds on a fixed tube over $I_2$. These bounds also hold at
$\eta=0$; none of the changes divides by $\eta$, $e$, or $n$.

To construct an invariant graph with common finite derivatives, embed
$I_2$ as an arc of a circle and extend the coefficients with a
fixed smooth cutoff equal to one on $I_1$. Extend the normal coefficient
$X$ to a function $\widetilde B\ge b_0>0$ on the circle tube.
The extended equations are
\[
 \dot\theta=e^2\widetilde A(\theta,w),\qquad
 \dot w=-\widetilde B(\theta,w)w
                 +q\widetilde R(\theta,w).
\]
The extension retains both small factors $e^2$ and $q$, and all its
coefficient norms are common over the prescribed compact parameter
sets. At $e=0$ the zero circle has zero tangential rate and a normal
contraction rate bounded below by $b_0$. For the fixed order $K+1$,
the normal-hyperbolicity inequalities persist after one parameter
shrinking. Fenichel's theorem
\cite[Theorem~3, p.~215]{Fenichel1971} therefore gives a $C^{K+1}$
invariant circle, whose projection is a graph $w=\phi(\theta)$.
It has common finite phase bounds: in the graph transform the normal
variations decay at least as $e^{-b_0t/2}$, while each fixed number
of base variations is bounded by a polynomial in $t$ times
$e^{C_Ke^2t}$. Taking $C_Ke^2<b_0/4$ gives common bounds for the
successive jet equations. The invariance equation then retains the
factor $q$, which is needed when the selected graph is translated back
to the physical variables.

That equation is
\begin{equation}\label{eq:small-ratio-graph-invariance}
 e^2\widetilde A(\theta,\phi)\phi'
       +\widetilde B(\theta,\phi)\phi
       =q\widetilde R(\theta,\phi).
\end{equation}
The differentiated-maximum induction in the proof of
Lemma~\ref{lem:pc-selected-graph} now retains $q$ on the right-hand
side of \eqref{eq:small-ratio-graph-invariance}: at order $m\le K$
the transport term vanishes at a maximum of $|\phi^{(m)}|$, the normal
coefficient is at least $b_0/2$, and every remaining term is $O_m(|q|)$
once the lower jets are bounded. The preliminary phase bounds control
the first-order quadratic term. Thus
$\|\phi\|_{C^K}\le C_K|q|$, uniformly also at $q=0$.
Restriction to
$I_0$ and reversal of the translation therefore give
\begin{equation}\label{eq:small-ratio-selected-graph}
 Y_s(X)=-\frac{e\eta}{X}+O_{C^K}(e^3|\eta|).
\end{equation}
In particular $\phi=0$ when $\eta=0$, so this estimate includes that
parameter face with the same constants.

Put $u=Y-Y_s(X)$ and
$\chi(X,u)=1+bnX-n^2e^2(2Y_s(X)+u)$. Subtracting the graph
invariance equation gives, in the original time,
\[
 X'=e^2\mathcal F(X,Y_s(X))+e^2\chi(X,u)u,\qquad
 u'=\{X-e^2Y_s'(X)\chi(X,u)\}u.
\]
The normal coefficient is a positive unit. The graph estimate supplies
common $C^{2j+10}$ bounds for these coefficients, so we may apply
Lemma~\ref{lem:uniform-normal-product} with $\tau=\rho=e^2$ at order
$J=j+2$. This gives exact product coordinates
\[
 \theta'=e^2f(\theta),\qquad \omega'=a(\theta)\omega,
 \quad f=1-\theta+Cc\theta^2+O_{C^1}(e|\eta|),\quad a\ge a_*>0.
\]
Their inverse satisfies $X_\omega=O_{C^j}(e^2)$ and
$Y=Y_s(\theta)+\omega u_1(\theta,\omega)$, with
$u_1$ and $X_\theta$ positive units. Since
$f_\theta=-1+O(c+e|\eta|)$, its root is simple. The scalar
linearization at the end of that lemma, with $g=f/a$ and
$\beta=-g'(\theta_s)>0$, gives \eqref{eq:simple-exact-power}
with $r=e^2\beta\asymp e^2$. The two extra product orders leave the
asserted $C^j$ bounds after this scalar coordinate change.

The estimate $X_\omega=O_{C^j}(e^2)$ also controls changes to physical
sections despite the weak longitudinal speed. Fix a
canonical base $\theta_b$ a distance comparable to $\sigma>0$ from
the simple root, and then a normal collar of width $\delta\ll\sigma$.
In the inverse product coordinates
$X=\theta+e^2\omega L(\theta,\omega)$, the longitudinal displacement
across this collar is $O(e^2\delta)$. A vertical section just to the
left of the incoming canonical section can therefore be chosen so that
the intervening forward flight has duration $O(\delta/\sigma)$.
For a nearby physical section $X=\theta_b+e^2\ell$, the hitting
equation, starting on $\theta=\theta_b$, becomes
\[
 \int_0^T f(\Theta(t))\,dt
       +\omega(T)L(\Theta(T),\omega(T))-\ell=0
\]
after division by $e^2$. Its time derivative is separated from zero
on the chosen collar because $|f(\theta_b)|\asymp\sigma$.
The implicit-function theorem therefore gives common derivatives
through the required order for this section map and its inverse.
The intervals, $\sigma$, and $\delta$ are fixed before shrinking
$e,c,|\eta|$; no relation between their rates of approach to zero is
used.

Finally put $\mathcal P=D\circ H$.  At a point where $\mathcal P'=1$,
\[
 \mathcal P''=\frac{D''}{D'^2}+\frac{H''}{H'}
              =\frac{r-1}{rD(H)}+\frac{H''}{H'}<0,
\]
since $D\le\xi_0$ and \eqref{eq:simple-small-ratio-curvature} holds.
Every zero of $\mathcal P'-1$ is a strict downward crossing.
The single-maximum argument of Lemma~\ref{lem:pc-ratio} therefore gives
at most two simple return zeros or one exactly quadratic zero on the
actual interval, with the multiplicity transfer of
Section~\ref{sec:return-framework}.
\end{proof}



\section{Simple-contact zero count}\label{sec:simple-contact}

We now compose the local passages of
Section~\ref{sec:simple-contact-tools} with the complete complementary
transition. Bounded saddle ratios use Lemma~\ref{lem:dulac-jets}; small
ratios use the exact power passage just constructed. Before estimating
zeros, we specify the return domain on which these compositions preserve
physical fixed points and their multiplicities.

\begin{proposition}\label{prop:simple-contact}
Suppose that $0<A\le A_+<\infty$ and that either
\begin{enumerate}
\item $C$ belongs to a compact subinterval of $(0,\infty)$; or
\item $A\ge A_0>0$, $C$ is bounded, and the polynomial satisfies
$P(z)\ge p_0>0$ on the fixed negative interval preceding the small
contact neighborhood used below.
\end{enumerate}
There are positive constants $c_0,d_0,\varepsilon_0,\rho_{\rm H}$
and a neighborhood $U$ of the whole graphic such that, for
$|c|<c_0$, $|d|<d_0$ and $0<\varepsilon<\varepsilon_0$,
the total actual finite minimal-return multiplicity of the isolated
periodic orbits $L\subset U$ satisfying
$d_{\rm H}(L,\Gamma_b)<\rho_{\rm H}$ is at most two.
The choices are locally uniform in
$b\in(0,2)$ and uniform on every compact set in the second case having
the stated positive lower bound.
\end{proposition}

\paragraph{Domain of the outgoing inverse.}
Let $D$ denote the local passage and $N$ the transition from the physical
source to its incoming coordinate. The conjugate return $D\circ N\circ C_+$
can be used only where the outgoing continuation $C_+$ and its inverse
are both defined.
Fix the original positive low section core $I_0$, the local outgoing
section $\Omega$, and a larger ordinary outgoing collar containing
$[0,a_{\rm out}]$. Write $I_{\rm out}=(0,a_{\rm out})$ for its counting interval.
The ordinary positive continuation $C_+:I_{\rm out}\to J\subset I_0$ is defined
on the entire collar, with $C_+'>0$, before the return is defined.
Reference barriers and the regular fast tube are constructed for that
whole collar and the same field. A smaller incoming collar puts every
finite local output in $I_{\rm out}$; there is no orbit-dependent collar choice.

The signed transition along the fast arc, descent and negative predecessor
is increasing in the physical source ordinate. Let $I_{\rm pass}$ be the
set of incoming ordinates admitting a positive finite contact passage
with $F>0$ to the target core. Between two passing
graphs, $F_y>0$ bounds the intermediate speeds below by the positive
minimum on the lower graph, for this fixed field. Scalar order then
continues every intermediate graph to the target; transverse finite
hits give openness. Thus $I_{\rm pass}$, and its source preimage
$I_{\rm ret}\subset I_0$, are open intervals, possibly empty.
The regular fast clearance, negative
predecessor and rightgoing contact event show that this is a minimal
first return $R$ on $I_{\rm ret}$.

For the central or compact passage, or the late weak-saddle class,
write $I_{\rm ret,local}$ for the actual local-return interval.
In the weak case it is the lower, late-class interval, excluding the
finite threshold assigned to the early class. The local entry and
outgoing-collar estimates give
\[
 R(I_{\rm ret,local})\subset J.
\]
This containment is asserted only for these local passages. The early
weak class retains its separate physical displacement and the signed
gluing argument below; no small weak-saddle collar is asserted for it.

Consequently $C_+^{-1}$ is used only on $I_{\rm ret,local}\cap J$:
\[
\begin{gathered}
 I_{\rm adm}=C_+^{-1}(I_{\rm ret,local}\cap J),\qquad
 H=N\circ C_+,\quad \mathcal P_a=D\circ H,\\
 R(C_+(a))=C_+(\mathcal P_a(a))\qquad(a\in I_{\rm adm}).
\end{gathered}
\]
Every point of this interval follows the same sequence of positive
finite passages. Flow uniqueness therefore gives the conjugacy on this
domain.
The conjugacy preserves actual minimal-return multiplicity by
Section~\ref{sec:return-framework}; its derivative need only be nonzero
for each fixed field.
Axis hits, zero amplitudes on
equilibria or stable separatrices, missing targets and infinite flights
are outside the actual return domain. The same domain requirements
apply to bounded signed $C$ under the second hypothesis of the proposition.

\subsection{Bounded ratios to the contact scale}\label{subsec:simple-bounded}

We first bound the rescaled constant term and forcing:
$|k|,|\ell|\le K$, where
$k=c/\varepsilon^2$ and $\ell=d/\varepsilon^3$.
Putting $n=A\varepsilon^2$, $z=nU$, $y=n^2V$ gives
\begin{equation}\label{eq:simple-bounded-system}
 U'=\mathcal F=V(1+bnU-n^2V)+k-U+C\varepsilon^2U^2,
 \qquad V'=UV+\ell.
\end{equation}
The tail estimate with $t=\varepsilon$ gives
$V(-L)\le L+K+1$.  On $[-L,0]$, above a sufficiently large fixed
level $T$, one has $\mathcal F\ge V/4$ and $V'\le K$.
Thus $(V^2)_U\le8K$ on each excursion above $T$, and
\begin{equation}\label{eq:simple-bounded-height}
 V(U)\le B_*:=\sqrt{T^2+8KL}\qquad(-L\le U\le0).
\end{equation}
Thus all incoming contact passages enter a bounded scaled strip. For
signed $C$, take $|C|\varepsilon^2L^2<1$ and increase $T$
by a fixed amount; the argument is unchanged.

\paragraph{Parameters separated from $\ell=0$.}
When the rescaled forcing is separated from zero, negative forcing
prevents passage and positive forcing makes the complete return
contracting. Fix $\ell_0>0$. If $\ell\le-\ell_0$, start at $U=-2L$, where
$V\le2L+K+1\le3L$.  On $[-2L,-L]$ a hypothetical passing orbit
has $V'\le-LV-\ell_0$ and $0<\mathcal F\le12L$, after increasing
the harmless numerical constant if necessary.  Consequently
\[
 V(-L)\le(3L+\ell_0/L)e^{-L/12}-\ell_0/L<0
\]
for a sufficiently large fixed $L$, a contradiction.

For $\ell\ge\ell_0$, the entrance instead lies below the positive
nullcline $m(U)=-\ell/U$.  To see this uniformly, on the first half
of $[-2L,-L]$ the level $2\ell/(3L)$ is an upper barrier and
\[
 V(-3L/2)\le2\ell/(3L)+3Le^{-L/16}<3\ell/(4L).
\]
If the orbit remained above $m$ on the second half, $V$ would be
nonincreasing, contradicting $m(-L)=\ell/L$.  At a meeting with $m$,
$d(V-m)/d\sigma=-m'\mathcal F<0$, so it cannot cross back upward.
It follows that $UV+\ell>0$ and
\begin{equation}\label{eq:simple-forcing-negative-density}
 \partial_V\frac{UV+\ell}{\mathcal F}
 =\frac{U\mathcal F-(UV+\ell)\mathcal F_V}{\mathcal F^2}<0
 \qquad(-L\le U<0).
\end{equation}
If there is an equilibrium at a negative abscissa, the first one,
$U_n$, prevents passage: the region
$\{-L\le U\le U_n,\ 0\le V\le m(U)\}$ is forward invariant.
Its upper boundary points inward, and on the right boundary
$\mathcal F(U_n,V)\le\mathcal F(U_n,m(U_n))=0$.
This includes the coalescence of two equilibria.

Every remaining crossing has $V(0)\ge v_0>0$ uniformly.  Indeed, on
$[-1,0]$, below $\min(1,\ell_0/2)$, one has
$V'\ge\ell_0/2$ and $0<\mathcal F\le K+4$.
Choosing $v_0$ smaller than both half that level and
$\ell_0/[4(K+4)]$, a terminal value below $v_0$ would force
$V(-1)<0$.  The level $v_0$ cannot be crossed downward there.
For $U\ge0$, a small parabola $V_b=v_0/2+\gamma U^2$ points inward:
\[
 V'-V_b'\mathcal F\ge\ell_0+
 U\{(1-4\gamma)v_0/2-2\gamma K\}
 +\gamma(1-4\gamma-2|C|\varepsilon^2)U^3>0.
\]
Take $U_1$ large enough that $\mathcal F\ge V/4$ for $U\ge U_1$.
On $[0,U_1]$ the bounded-height argument used in
\eqref{eq:simple-bounded-height} gives $V\le B_1$; since
$V'\ge\ell_0$, its flight time is at most $B_1/\ell_0$.
The vertical-section derivative on this compact part is
\[
 \frac{\mathcal F(0,V(0))}{\mathcal F(U_1,V(U_1))}
     \exp\!\int\operatorname{div}(\mathcal F,UV+\ell)
 \le K_1.
\]
No inverse of the possibly small speed at $U=0$ is used.
On the rest of the positive part, \eqref{eq:simple-Ry} and the lower
parabola give a logarithmic derivative loss at most
$K_2+K_2\varepsilon^2\log(1/(A\varepsilon^2))$.
The negative transition has derivative at most
$K_3\varepsilon\exp[-(t_0+t_1\log(A_+/A))/\varepsilon^2]$,
and \eqref{eq:simple-forcing-negative-density} adds no loss.
The derivative of the complete return is therefore bounded above by
\begin{equation}\label{eq:simple-forcing-multiplier}
 K_4\varepsilon\exp\left\{
 -\frac{t_0+t_1\log(A_+/A)}{\varepsilon^2}
 +K_4\varepsilon^2\log\frac1{A\varepsilon^2}\right\}<1.
\end{equation}
Choose $\varepsilon$ so that $K_4\varepsilon^4<t_1/2$; then the
coefficient of $\log(A_+/A)$ is negative and the remaining logarithms
are absorbed by $t_0/(2\varepsilon^2)$. The displacement is therefore
strictly decreasing on its single return interval and has at most one zero.

\paragraph{Small forcing and negative $k$.}
For small forcing, a negative constant term is excluded by reducing
the incoming height below the level needed for positive speed.
Fix $k_0>0$, let $k\le-k_0$, and take $|\ell|\le k_0^2/64$.
Put $H=k_0/8$.  On $[-2L,-L]$, as long as $V\ge H$,
\[
 V'=UV+\ell\le-LV+k_0^2/64\le-(L/2)V
\]
for fixed large $L$.  A passing orbit has $0<\mathcal F\le B L$ and
enters this strip below $3L$.  Until its first hit of $V=H$ it therefore
satisfies
\[
 \frac{dV}{dU}\le-\frac{V}{2B},\qquad
 V(-L)\le3Le^{-L/(2B)}<H.
\]
If it hits $H$ earlier, the vector field on that boundary points down and
prevents a later upcrossing.  Thus in all cases $V(-L)<H$.  For
$U\le-k_0/4$,
$V'|_{V=H}\le-k_0H/4+k_0^2/64<0$.
At $U=-k_0/4$ this implies
$\mathcal F\le2H-k_0+k_0/4+O(\varepsilon^2)<0$.
Thus this parameter range has no passing orbit.

\paragraph{The saddle-node at $(k,\ell)=(0,0)$.}
At these parameters the origin has eigenvalues $-1,0$ and center
equation $V'=V^2+O(V^3)$; $V=0$ is its strong invariant line.
Lemma~\ref{lem:saddle-node-passage} applies, with the center coordinate
oriented by $s_V(0)>0$. To use its small-passage estimate, we must first
put every physical entry in a common small incoming range.
For a small fixed $d_0>0$, descent on $[-2L,-L]$ and the constant
upper barrier on $[-L,-d_0]$ give
\begin{equation}\label{eq:simple-central-entry}
 V(-d_0)\le
 \max\{(2L+K+1)e^{-k_5L},\,2|\ell|/d_0\}.
\end{equation}
For $|k|<d_0/4$, the speed on the latter interval is at least $d_0/2$.
Choose $L$ large, then $|k|,|\ell|$ small.  In product coordinates
the entry has $w\in[3d_0/4,5d_0/4]$ and
$|s-\sigma|\le K_5(V+|k|+|\ell|)$.
It reaches $w_0=d_0/2$ in the bounded time $\log(w/w_0)$, with a
positive transverse derivative and bounded derivatives of the section
change.  At zero parameters, positive physical entries have $s>0$;
their potentially infinite time from the right is included in
\eqref{eq:simple-saddle-node-smallness}.

For the positive reference, use the linear and parabolic barriers in
the saddle-node part of the proof of
Lemma~\ref{lem:lower-positive-references}, with
$p(U)=k-U+C\varepsilon^2U^2$, $\ell_*=\ell$ and
$n=A\varepsilon^2$. Start at $U=\delta$ on
$h=\delta+\delta^2+O(\delta^3+|k|+|\ell|+\varepsilon^2)$.
The lower line $(1+\delta/4)U$ has a strict crossing margin for
$|k|+|\ell|\ll\delta^2$; the small and large parabolas then give
$h\asymp1+U^2$ and $\mathcal F/h\ge k_6>0$.
These inequalities do not use a lower bound for $A$: their physical
smallness is $y\le M(n^2+z^2)$, with $n\le A_+\varepsilon^2$.
They also allow signed bounded $C$, since its adverse contribution is
$O(\varepsilon^2h)$. Fix $z_c$ before shrinking $n$.
Thus Lemma~\ref{lem:positive-transition} applies uniformly in $A$.

Let $N(u)$ be the complete transition from the positive section to the
incoming center coordinate.  The convex map of
Lemma~\ref{lem:negative-transition} remains convex up to $U=-d_0$:
in \eqref{eq:simple-Ryy} the positive term $xp(1+bz)$ dominates
the forcing, since $|\ell|\ll d_0^2$, and all higher physical-height
terms tend uniformly to zero on the bounded core.
Its scaled first derivative is $o(1)$ by
\eqref{eq:simple-scale-cancellation}.  The final bounded section
change therefore gives $N'>0$ and $N''/N'\ge-o(1)$.
Combining this incoming estimate with the outgoing transition gives,
for $H=N\circ C_+$,
\begin{equation}\label{eq:simple-central-curvature}
 aH''/H'=aC_+''/C_+'+aC_+'(N''/N')\ge-K_7a.
\end{equation}
First choose the product neighborhood with $W_s\le1/4$, then
$a_0$ with $K_7a_0\le1/4$.  Decrease the incoming range and the
parameters in \eqref{eq:simple-central-entry} so that
\eqref{eq:simple-saddle-node-smallness} puts every fixed point in
$0<a<a_0$.  The actual passing condition restricts $H$ to one component
to the right of any scalar root, hence to an interval.  Intersecting
with $(0,a_0)$ keeps an interval containing every fixed point.
Equations \eqref{eq:simple-central-curvature} and
\eqref{eq:simple-pfaff-hypothesis} give at most two cycles.

\paragraph{Small forcing and positive $k$.}
For $k\in[k_0,K]$ the saddle near $(U,V)=(k,0)$ persists, with
\[
 U_s=k+O(\varepsilon^2+|\ell|),\quad V_s=-\ell/U_s,
 \quad\lambda_-=-1+O(\varepsilon^2+|\ell|),\quad
 \lambda_+=k+O(\varepsilon^2+|\ell|).
\]
Its characteristic ratio $r=-\lambda_-/\lambda_+$ lies in a compact
subinterval of $(0,\infty)$, including values on either side of one.
The descent \eqref{eq:simple-central-entry}, followed to
$U=U_s-\delta$, makes the incoming unstable coordinate $s$ as small
as prescribed.  That continuation has $\mathcal F\ge\delta/2$;
its derivatives are uniformly bounded.
Choose the saddle coordinates with $U_x>k_8>0$, $U_y<-k_8$.
In the prepared form used in Lemma~\ref{lem:dulac-jets}, a negative
entry $s$ exits on the left with stable coordinate at most
$y_0(|s|/x_0)^{r_-/2}$.  For small $|s|$, its physical longitudinal
speed is negative there.  Such a trajectory, and the infinite-time
entry $s=0$, cannot be a counted contact passage.  Hence $s>0$.

At $U=U_s+\delta$ the reference satisfies
$h=(k+1)\delta+O(\delta^2+|\ell|+\varepsilon^2)$.
Use the bounded-ratio construction in the proof of
Lemma~\ref{lem:lower-positive-references}, replacing $k_c$ by $k_0$
and $\ell_*$ by $\ell$. The lower line
$(1+k_0/2)(U-U_s)$ has a strict positive crossing margin; the same
parabolas give $h\asymp1+U^2$ and $\mathcal F/h\ge k_9>0$.
Again only $n=A\varepsilon^2\le A_+\varepsilon^2$ enters the physical
smallness, and bounded signed $C$ contributes $O(\varepsilon^2h)$.
The positive transverse orientation and compact local section change
give the incidence required by Lemma~\ref{lem:positive-transition}.
Hence \eqref{eq:simple-positive-jets} holds with $n=A\varepsilon^2$.

For the negative transition to the saddle's incoming section,
Lemma~\ref{lem:negative-transition} and its final bounded section
change give $N'>0$, $N'=o(1)$ and $|N''/N'|\le K_9\varepsilon^{-2}$.
Consequently
\begin{equation}\label{eq:simple-compact-complement}
 0<H'\le m(\varepsilon)\longrightarrow0,\qquad
 |H''/H'|\le K_{10}+K_{10}n^{3/2}\varepsilon^{-2}
             \le K_{11},\qquad H=N\circ C_+.
\end{equation}
The identity $n^{3/2}\varepsilon^{-2}=A^{3/2}\varepsilon$ is what
makes the bound uniform in $A$.

The smallness of $H'$ allows the two-zero comparison to include $r=1$.
At a point where $(D\circ H)'=1$, set
$s=H(a)$, $T=\log(1/s)$ and $M=\log(1/m)-\log K_{12}$.
Lemma~\ref{lem:dulac-jets} gives
$(1-r)T\ge M$.  For large $M$ this is impossible if $r\ge1$.
Otherwise $\beta=1-r\ge M/T$ and
$sD''/D'\le-\beta/2$, since $T e^{-\nu T}$ is bounded.
At that point
\[
 (D\circ H)''=\frac{D''}{D'^2}+\frac{H''}{H'}
 \le-\frac{\beta}{2sD'}+K_{11}<0,\qquad
 \frac{\beta}{sD'}\ge\frac{M e^{r_-T}}{K_{12}T}.
\]
Taking $m$ small proves the last strict inequality.
All zeros of $(D\circ H)'-1$ are downward crossings, so there is at
most one. A displacement zero at this stationary point is exactly
quadratic and the sole zero; all other zeros are simple. Hence the
total actual fixed-point multiplicity is at most two.
Fix the outgoing collar first; choose the incoming collar small enough
that $D$ maps into it and that the negative-entry exclusion holds;
then choose $L,|\ell|,\varepsilon$.  Monotonicity makes the resulting
return domain one interval.  This order places every counted cycle
in the domain on which the estimates were proved.

These cases cover every fixed bounded $(k,\ell)$ set with common
thresholds. Choose a small central neighborhood, then $k_*>0$
inside it, and finally $\ell_*>0$ below its threshold, $k_*^2/64$,
and the positive-saddle threshold.  The five parameter ranges
\[
 \ell\ge\ell_*,\quad \ell\le-\ell_*,\quad
 |\ell|<\ell_*, k\le-k_*,\quad
 |\ell|<\ell_*, |k|<k_*,\quad
 |\ell|<\ell_*, k\ge k_*
\]
have bounds $1,0,0,2,2$, respectively.  Equalities are assigned to a
range whose estimate holds with a strict margin.

\subsection{Ratios larger than the contact scale}\label{subsec:simple-radial}

Put
\[
 \rho=\max\{\varepsilon,\sqrt{|c|},|d|^{1/3}\},\qquad
 e=\varepsilon/\rho,\quad \kappa=c/\rho^2,\quad\eta=d/\rho^3.
\]
The compact direction set is
\begin{equation}\label{eq:simple-direction-set}
 0\le e\le1,\quad |\kappa|,|\eta|\le1,\qquad
 \max\{e,\sqrt{|\kappa|},|\eta|^{1/3}\}=1.
\end{equation}
For $e$ bounded below, $k=\kappa/e^2$ and $\ell=\eta/e^3$ are bounded,
so the preceding subsection applies. It remains to treat small $e>0$,
when the constant term or forcing determines the larger scale $\rho$.

\paragraph{Negative forcing or a negative constant term.}
On the larger scale, the same incoming-height argument excludes
negative forcing and a negative constant term. Use the exact scaling
\[
 X=\frac{z}{A\rho^2},\quad
 Y=\frac{y}{A^2\varepsilon^2\rho^2},\quad
 X'=e^2\mathcal F,\quad Y'=XY+e\eta,
\]
where
\begin{equation}\label{eq:simple-radial-system}
 \mathcal F=YQ+\kappa-X+C\rho^2X^2,\qquad
 Q=1+bA\rho^2X-A^2\rho^4e^2Y.
\end{equation}
Lemma~\ref{lem:negative-transition}, with $t=\rho$, gives
$Y(-L)\le L+K+1=:B_0$.
If $\eta\le-\eta_0<0$, on the unit interval $[-L,-L+1]$
one has $Y'\le-e\eta_0$, $Y\le B_0$ and
$0<\mathcal F\le B_1$.  Thus
$Y_X\le-\eta_0/(eB_1)$, which is incompatible with $Y>0$ when
$e$ is small.

If $\kappa\le-\kappa_0<0$, take the section $X_g=\kappa/2$.
Its polynomial value is at most $-\kappa_0/4$ after decreasing $\rho$.
Choose $\theta>0$ with $2\theta<\kappa_0/8$.
On the same first unit interval, as long as $Y\ge\theta$,
$Y'\le-(L-1)Y/2$; hence
\[
 Y(-L+1)\le B_0\exp[-(L-1)/(2e^2B_1)]<\theta
\]
unless it already met that level.
For $X\le-\kappa_0/2$, the boundary $Y=\theta$ still points down
when $e|\eta|<\kappa_0\theta/4$.
It therefore reaches $X_g$ with
$\mathcal F\le2\theta-\kappa_0/4<0$, another contradiction.
Both exclusions use the incoming-height bound from the fixed physical
section; that bound is essential to the scaled argument.

\paragraph{Positive forcing separated from zero.}
Assume $\eta\ge\eta_0>0$. To retain contraction between the negative
tail and the core, we first place the physical incoming orbit below
the positive nullcline.
On the far negative tail set
$h_s(z)=-\Delta/z$, $w=y-h_s$ and
\[
 E(z,s)=\exp\!\int_s^z\frac{\zeta}{F(\zeta,y(\zeta))}\,d\zeta.
\]
The exact difference equation and its integral are
\begin{equation}\label{eq:simple-radial-incidence}
 w_z=(z/F)w-h_s',\qquad
 w(z_g)=E(z_g,-q_g)w(-q_g)
                 -\int_{-q_g}^{z_g}E(z_g,s)h_s'(s)\,ds,
 \quad z_g=-LA\rho^2 .
\end{equation}
The first term has absolute value at most
$K\varepsilon^3 e^{-\gamma T_A/\varepsilon^2}$.
On the final interval of length $A\varepsilon^2$,
the kernel is bounded below by a positive constant and
$h_s'\ge k\Delta/(A^2\rho^4)$.
The integral is consequently at least
$k A^2\varepsilon^5\eta/\rho$.
The ratio of the first term to this lower bound is at most
\[
 K\frac{\rho}{\varepsilon^2}A^{-2}
    \left(\frac{A(1+L\rho^2)}{q_g+A}\right)^{\gamma/\varepsilon^2}
       \longrightarrow0.
\]
The maximization in $A$ from Lemma~\ref{lem:negative-transition}
makes this limit uniform for arbitrarily small relative scales.
Thus $w(z_g)<0$. At $w=0$ its derivative is
$-h_s'<0$, so the passing orbit remains below the nullcline up to zero.
In particular $G>0$ there, and
$R_y=(zF-GF_y)/F^2<0$ on the entire negative contact interval.

For the positive part, the forcing determines the scale
$s=A\varepsilon\rho$. Put
$z=sU$, $y=s^2V$.  The equations are
\begin{equation}\label{eq:simple-radial-core}
 U'=V(1+bsU-s^2V)+\kappa-eU+C\varepsilon^2U^2,\qquad
 V'=UV+\eta.
\end{equation}
The low-height argument on $[-1,0]$ used above gives
$V(0)\ge2v_0>0$.  With $\gamma>0$ small,
$V_b=v_0+\gamma U^2$ is a strict lower barrier: its crossing
expression is bounded below by
$\eta_0+U(v_0/2+\gamma U^2/2)>0$.
After excluding the already treated parameters
$\kappa\le-\kappa_*<0$, choose
$\kappa_*\le v_0/8$ and $e\le\sqrt{v_0\gamma}/4$.
Then $U'\ge V/4$ throughout the positive part.
For signed $C$, absorb
$|C|\varepsilon^2U^2\le(|C|\varepsilon^2/\gamma)V$
and replace $V/4$ by $V/8$.
The estimate
$y_z\le Kz+K s^3/(s^2v_0+\gamma z^2)$ keeps a sufficiently small
contact height in the physical neighborhood up to $z_c$.

In the scalar variational numerator of \eqref{eq:simple-radial-core},
the terms $-eU^2$ and $-\eta(1+bsU-s^2V)$ are nonpositive.
The other terms, divided by the square of the speed, have integrals
bounded by
\[
 K\int_0^{z_c/s}\frac{U\,dU}{(v_0+\gamma U^2)^2}
 +K\varepsilon^2\int_0^{z_c/s}
                         \frac{U^3\,dU}{(v_0+\gamma U^2)^2}
 +Ks^2\int_0^{z_c/s}\left(U+\frac1{v_0+\gamma U^2}\right)dU.
\]
This is at most $K+K\varepsilon^2\log(1+z_c/s)$.
Combining it with the negative estimate gives
\[
 \log\mathcal P'
 \le K-\frac{\gamma}{\varepsilon^2}
             (t_0+t_1\log(A_+/A))
       +K\varepsilon^2\{\log(A_+/A)+2|\log\varepsilon|+K\}<0.
\]
Here $\rho\ge\varepsilon$ was used in the last logarithm.
All scale factors cancel in vertical-section derivatives.
Thus the complete return has at most one fixed point on this
parameter range.

\paragraph{A positive constant term and small forcing.}
The last case has $\kappa=1$ and $|\eta|$ small.
Write $c=\rho^2>0$, $n=Ac$, $e=\varepsilon/\sqrt c$.
The equations are \eqref{eq:simple-small-ratio-system}.
Orbits may leave a small neighborhood of $Y=0$ before reaching the
saddle, so a local saddle map alone does not cover the return domain.
We separate these early departures from the late passages near the saddle.
Fix a small $\sigma>0$, then a smaller $\delta>0$.
The separating event is the first upward intersection with $Y=\delta$:
an intersection at $X_b\le1-\sigma$ defines the first class.
The second, late class enters a saddle neighborhood near $X=1-2\sigma$
below $Y=\delta$; its fixed outgoing normal section has height between
$3\delta$ and $5\delta$.

We first show that these classes exhaust the incoming orbits and form
ordered source intervals.
For $\eta\le0$ the homogeneous negative estimate gives
\begin{equation}\label{eq:simple-small-forcing-entry}
 Y(-L)\le\frac{K e}{A^2\sqrt c}
        \exp\left(-\frac{\gamma T_A}{ce^2}\right)
                  \le e^{-B/e^2}
\end{equation}
for any prescribed $B>0$, by taking the upper bound for $c$ small first.
The factor $A^{-2}$ is absorbed by maximizing as before, and the
remaining $c^{-1/2}$ by $\exp[-k/(ce^2)]$.
On $[-L,0]$ the ordinate decreases; on $[0,1-\sigma]$ it satisfies
$Y_X\le K Y/(\sigma e^2)$.
Taking $B>K/\sigma$ proves that this sign of forcing cannot belong
to the first class.
For $\eta\ge0$, the height on $[-L,0]$ is uniformly bounded.
For example an upward crossing of a large fixed level $T$ requires
$X\ge-e\eta/T$; while $Y\ge T$, one has
$\mathcal F\ge Y/2$ and $(Y^2)_X\le4\eta/e$.
The remaining interval has length at most $e\eta/T$.
This bounds the height and gives, by an integrating factor,
\[
 Y(0)\le B_0e^{-L^2/(K_1e^2)}
                +\frac{K_1\eta}{e}
                     \int_{-L}^0e^{-X^2/(K_1e^2)}\,dX
       \le B_0e^{-L^2/(K_1e^2)}+K_2\eta<\delta/4.
\]
An upward crossing of $Y=\delta$ at negative $X$ could not be
followed by a downward crossing before zero, since
$X\delta+e\eta$ increases with $X$.
Its first upward intersection therefore has $X_b>0$.
Every counted cycle eventually has such an intersection because its
outgoing physical height is separated from the shrinking level
$n^2e^2\delta$. Write $q$ for the increasing physical source ordinate.
Scalar ordering gives $Y_q>0$ at every fixed pre-hit abscissa, and at
the upward event $Y_X>0$. Thus
\[
 (X_b)_q=-Y_q/Y_X<0.
\]
The second class is therefore the lower source interval and the first
class the upper interval, separated, when both occur, by one finite
transverse threshold.

For the early class, $\eta\ge0$. We estimate the complete multiplier
by following the growing height through the core. Use $Y$ as independent
variable up to $Y=e^{-2}$. A first zero of $\mathcal F$ is impossible, because
its derivative there is $\mathcal F_Y>0$.
Moreover,
\begin{equation}\label{eq:simple-growing-height}
 \frac{d(X^2)}{dY}\le K_3e^2+\frac{K_3e^2}{Y}
                         +\frac{K_3ce^2X^2}{Y}.
\end{equation}
Gronwall's inequality keeps $X$ bounded.  Before a fixed large level
$Y_0$, its squared increment is $O(e^2)$, leaving
$X\le1-\sigma/2$.  Choose $Y_0$ above four times the resulting
global upper bound for $X$; above it $\mathcal F\ge Y/4$.
Thus $\mathcal F\ge k_1(1+Y)$ throughout.
The terminal $X$ is bounded below as well.  If $XY<e\eta$ above
$Y_0$, then $X_Y\ge eY/(8\eta)$, which forces $XY=e\eta$ before
$Y_0+1$, for sufficiently small $\eta$.  Thereafter
$(X^2)_Y\ge e^2/4$.  For $\eta=0$ this last estimate applies directly.
It follows that at $Y=e^{-2}$,
\[
 0<x_-\le X_t\le x_+<\infty,\qquad y=n^2,\quad z=nX_t.
\]
On the whole scaled positive part the exact density is
\[
 \partial_Y\frac{XY+e\eta}{e^2\mathcal F}
 =\frac{X[1-X+CcX^2+n^2e^2Y^2]
                 -e\eta(1+bnX-2n^2e^2Y)}
        {e^2\mathcal F^2}.
\]
Its integral is at most $K_4/e^2$.
From the terminal point a small lower parabola and a large upper
parabola in $U=z/n,V=y/n^2$ give
$y\asymp n^2+z^2$ and $F\ge k_2y$ to the fixed positive section.
At a lower meeting the term $UV(1-2\gamma Q)$ is positive; the terms
of orders $e^2U$ and $\varepsilon^2U^3$ are absorbed after fixing
$\gamma$. After shrinking the physical neighborhood, at the upper
parabola $V=M(1+U^2)$ one has $F/n^2\ge V/2$, or equivalently
$F\ge y/2$, and $dV/dU<2MU$.
The remaining density integral is therefore at most
$K_5+K_5\varepsilon^2\log(1/n)$, by the integration in
Lemma~\ref{lem:positive-transition}.
Since the negative part has $R_y<0$, the complete derivative satisfies
\begin{equation}\label{eq:simple-first-class-multiplier}
 \log\mathcal P'\le
 \log(K_6\varepsilon)-\frac{\gamma T_A}{ce^2}
                 +\frac{K_6}{e^2}+K_6
                 +K_6\varepsilon^2\log(1/n)<0.
\end{equation}
Indeed $\log(1/n)\le T_A+K_7+\log(1/c)$.
First take $c$ small to absorb $K_6/e^2$, then $\varepsilon$ small
to absorb $\varepsilon^2T_A$; the remainder
$ce^2\log(1/c)$ tends uniformly to zero.
There is at most one cycle in the first interval.

For the late class, we apply Lemma~\ref{lem:small-ratio-saddle}
after verifying the physical amplitude and the derivatives of both
complementary transitions. The incoming amplitude must be positive.
For $\eta\ge0$, the same selected weak graph is
\[
 Y_s(X)=-e\eta/X+O_{C^K}(e^3\eta).
\]
It is below the axis for $\eta>0$ and equals the axis for $\eta=0$.
The product inverse preserves the sign of $Y-Y_s$; hence
$\omega\le0$ is nonphysical in $Y>0$. For $\eta<0$, zero amplitude
approaches the saddle in infinite time, while a negative amplitude,
if initially physical, reaches the axis while growing in modulus.
The axis-hit argument is used only for this negative sign of the forcing.
The incoming tilted section is reached with bounded derivatives by
the explicit $O(e^2\delta)$ section adjustment in that lemma.
The transition from $-L$ to that section stays to the left of
$1-\sigma$ with bounded height and $\mathcal F\ge k_\sigma>0$.
Writing this scalar transition as $Q_e$ and its equation as
$Y_X=f_e(X,Y)$, direct differentiation on the fixed strip gives
\[
 |f_{e,Y}|+|f_{e,YY}|\le K_8e^{-2},\qquad
 0<Q_e'\le e^{K_8/e^2},\qquad
 \left|\frac{Q_e''}{Q_e'}\right|
       \le K_8e^{-2}e^{K_8/e^2}.
\]
The rescaled incoming map $v$ supplied by
\eqref{eq:simple-negative-jets} has
\[
 v'\le L_0(A,c,e)e^{-\gamma T_A/(ce^2)},\qquad
 |v''/v'|\le K_8\varepsilon^{-2}.
\]
Here $L_0$ has at most polynomial growth in the inverse small scales.
After taking $c$ small, the negative action absorbs the factor
$e^{K_8/e^2}$ and every polynomial loss, uniformly as $A\downarrow0$.
Consequently $(Q_e\circ v)'=o(1)$ and
\[
 \frac{(Q_e\circ v)''}{(Q_e\circ v)'}
 =\frac{v''}{v'}+\left(\frac{Q_e''}{Q_e'}\circ v\right)v'
 =O(\varepsilon^{-2}).
\]
Including the ordinate rescaling and the regular tilted-section change
therefore gives
\begin{equation}\label{eq:simple-small-ratio-negative}
 N'>0,\qquad |N''/N'|\le K_9\varepsilon^{-2}.
\end{equation}

To bound the outgoing transition independently of the weak eigenvalue,
move the normal outlet to a horizontal section
$Y=\delta_1\asymp\delta$, slightly above it.  Its $X$-coordinate is
a decreasing smooth function of the outgoing weak coordinate $a$,
with derivative a negative unit.
The inverse-coordinate bounds in
Lemma~\ref{lem:small-ratio-saddle} give uniform first three derivatives
for this conversion; shrink $e|\eta|\ll\delta$ after fixing $\delta$.
Continue with $Y$ as independent variable to $e^{-2}$.
Here $X$ stays in a positive compact interval,
$XY+e\eta\ge k_3XY$, and for $j\le3$,
\[
 \left|\partial_X^j
 \frac{e^2\mathcal F}{XY+e\eta}\right|
                         \le K_{10}e^2(1+1/Y).
\]
Its integral is bounded by
$K_{11}+K_{11}e^2|\log\delta|$.
After the last shrinking of $e$ this gives a nonzero first derivative
and bounded second and third derivative ratios to the terminal
horizontal section.
The parabolic bounds just established, a fixed compact transition to
a vertical $U$-section, and
Lemma~\ref{lem:positive-transition} give
$0<C_+'\le K_{12}n^{3/2}$ and
$|C_+''/C_+'|\le K_{12}$.
The coordinate $a=0$ is the actual outgoing unstable separatrix.
Its fixed small collar can be chosen by first fixing $\sigma$, then
$\delta\ll\sigma$, then the parameters: the horizontal-section
conversion has bounded derivative ratios independently of the last
choice of $\delta$, since its two heights have a bounded ratio.
This avoids a dependence of the allowable collar on an inverse weak
speed.

For the complete complementary transition $H=N\circ C_+$,
\begin{equation}\label{eq:simple-weighted-curvature}
 r|H''/H'|
 \le K_{13}e^2+
      K_{13}\frac{e^2n^{3/2}}{\varepsilon^2}
 =K_{13}e^2+K_{13}A^{3/2}\sqrt c\longrightarrow0.
\end{equation}
The weight $r\asymp e^2$ cancels the loss from the incoming second
derivative; the available unweighted estimate need not tend to zero.
Lemma~\ref{lem:small-ratio-saddle} gives at most two cycles in the
second interval.
The equality $X_b=1-\sigma$ is assigned to the first interval; its
transverse event and all estimates extend there, so it adds no cycle.

\begin{lemma}[Gluing the two simple-contact classes]
\label{lem:simple-two-class-gluing}
The two classes just constructed have total actual minimal-return
multiplicity at most two, including a finite threshold only once.
\end{lemma}
\begin{proof}
For $\eta\le0$ the first class is absent, so only the second-class
bound is needed. For $\eta>0$, let $\mathcal P(q)$ be the complete
return on the increasing physical ordinate and put
$\mathcal F(q)=\mathcal P(q)-q$. The two class descriptions are
restrictions of this same finite flow return. Its physical target is
transverse, so its hitting-time formula gives a $C^1$ return through
the finite threshold. The actual passing domain is one interval by
scalar order and $F_y>0$, as in Lemma~\ref{lem:sections}.
On the upper interval, \eqref{eq:simple-first-class-multiplier} gives
$\mathcal F'<0$, including the threshold assigned to that class.

On the lower, late interval use the increasing physical continuation
$q=C_+(a)$ only on $I_{\rm adm}=C_+^{-1}(I_{\rm ret,local}\cap J)$, where
$J=C_+(I_{\rm out})$. The complete outgoing collar was fixed first and contains
all late-class local outputs; its inverse is not used outside $J$.
The amplitude return is
$\mathcal P_a(a)=D(H(a))$, with $H=N\circ C_+$; the chosen
outgoing collar contains its outputs. The proof of
Lemma~\ref{lem:small-ratio-saddle} shows that every zero of
$\mathcal P_a'-1$ has $\mathcal P_a''<0$. Hence the amplitude displacement
has at most one stationary point, a strict maximum. If it has two
zeros, they cross first upward and then downward; a tangent zero is
a strict maximum and is the sole zero in that interval. Since
\[
 \mathcal F(C_+(a))=C_+(\mathcal P_a(a))-C_+(a),
\]
the increasing coordinate change preserves these signs.
Writing $J=(j_-,j_+)$, the containment
$\mathcal P(I_{\rm ret,local})\subset J$ also gives
$\mathcal F(q)>0$ for $q\le j_-$ and
$\mathcal F(q)<0$ for $q\ge j_+$ wherever the late return is defined.
Thus these signs extend across the parts of the late interval outside $J$.

Every lower zero is simple unless it is the sole, exactly quadratic
strict-maximum zero. If there are two simple lower zeros, or one
quadratic lower zero, the physical displacement is negative after the
last lower zero and up to the threshold. A zero at that threshold
would have negative derivative and therefore positive displacement
immediately to its left, a contradiction. The strictly decreasing
upper displacement has no zero either. Otherwise there is at most one
simple lower zero and at most one simple upper zero. Thus the sum of
actual multiplicities is at most two. The same
strict derivative statements exclude intervals of zeros. A finite
threshold zero is simple because $\mathcal F'<0$ there and is counted
once in this total multiplicity count; zero-amplitude separatrices and
infinite-time passage boundaries are not periodic orbits.
\end{proof}

All these arguments remain valid in the signed-$C$ case of the
proposition.  On fixed scaled intervals replace
$1-X+CcX^2\ge1-X$ by half that bound, after fixing the intervals and
decreasing $c$.  Near the saddle its derivative is still
$-1+O(c+e|\eta|)$.  On the positive parabolas the only new adverse
term is $-|C|\varepsilon^2U^2=O(\varepsilon^2V)$.
The negative estimates were already proved with
$P(-x)\asymp Ax$, and $A\ge A_0$ only improves their scale
cancellation.  These explicit replacements cover every sign-sensitive
inequality above.

\begin{proof}[Proof of Proposition~\ref{prop:simple-contact}]
Choose a small fixed $\eta_*>0$ for the positive-constant,
small-forcing argument, then $e_*>0$ small enough for it and both
separated-forcing arguments.
In \eqref{eq:simple-direction-set}, $e\ge e_*$ gives a fixed bounded
$(k,\ell)$ set and hence the bound two from
Subsection~\ref{subsec:simple-bounded}.
For $e<e_*$, $\eta\ge\eta_*$ gives the bound one and
$\eta\le-\eta_*$ gives zero.
In the remaining range $|\eta|<\eta_*$, compact normalization forces
$|\kappa|=1$: $\kappa=-1$ is excluded and $\kappa=1$ gives the bound
two by Lemma~\ref{lem:simple-two-class-gluing}. The parameter thresholds
are fixed before taking the common small upper bound on $\rho$.
Only finitely many choices are needed.
Within each range, fixed section sizes, barrier levels, and collar
widths are selected first.  The common upper bounds for $L^{-1}$ and
the small parameters are then chosen in the order used in the
corresponding proof, so estimates such as
$n^{2-K\varepsilon^2}\le n^{3/2}$ hold simultaneously on the whole
compact parameter set.

For each fixed field the subdivision is made by its coefficients,
so bounds from different parameter ranges are not added.  Within the
last range the two first-intersection intervals are the only
subdivision of the physical source.
To verify the interval assertions used throughout, two passing scalar
graphs bound every intermediate graph by uniqueness.  Since $F_y>0$,
the lower graph supplies positivity of $F$ and $y$, and the upper graph
prevents an earlier upper exit.  Intermediate graphs therefore reach
the same section.  Conditions on an incoming amplitude or on the
ordered first intersection are inverse images of intervals under
monotone maps.  Rolle's theorem has always been applied on these
intervals, not across missing portions of a return map.
By Lemma~\ref{lem:sections}, one fixed physical source contains the
single intersection of every sufficiently close whole-graphic cycle.
Thus no additional sum over source sections is required, and the
proposition follows.
\end{proof}
